% GENERATED FILE. Edit canonical lessons, then run npm run generate:books.
% canonical-source-hash: fnv1a64:7bd9870221157aa6
% canonical-lessons: JA-C141-raishuu, JA-C141-raigetsu, JA-C141-rainen, JA-C141-tsumori, JA-C141-deshou

\chapter{Next Week, Next Year: Plans and the Future}
\label{ch:ja-next-week-next-year-plans-and-the-future}
\hlchaptermodality{141}

\begin{chapteropening}
\textbf{By the end of this chapter:} \emph{I can say what I will do, what I plan to do and what will probably happen, and when: next week, next month, next year.}

Say what you will do next week, next month and next year with a verb in -masu, give one plan with tsumori desu, and make one guess with deshou.
\end{chapteropening}

\section[\texorpdfstring{\ja{らいしゅう} (raishū)}{らいしゅう (raishū)}]{\ja{らいしゅう} (raishū) — next week}

\label{lesson:JA-C141-raishuu}



\begin{quote}
Four recalls before the new word.

\begin{itemize}
\raggedright
  \item \emph{Recall:} say \emph{it is over}, politely — \textbf{R1}, one lesson back
  \item \emph{Recall:} say \emph{last week} — \textbf{R2}, five lessons back
  \item \emph{Recall:} say \emph{tempura} — \textbf{R3}, twenty lessons back
  \item \emph{Recall:} say \emph{outside} — \textbf{R4}, eighty lessons back
\end{itemize}
\end{quote}

\subsection*{You'll want to know: \ja{らいしゅう}}
\textbf{\ja{らいしゅう}} — \emph{raishū} — \textquotedblleft{}next week\textquotedblright{}: \textbf{\ja{せんしゅう}}, \textbf{\ja{こんしゅう}}, \textbf{\ja{らいしゅう}}, last week, this week, next week.

Japanese has no separate future tense. The polite present \textbf{\ja{ます}} says what will happen too, and a time word places it:

\begin{quote}
\textbf{\ja{らいしゅう} \ja{はたらきます。}} — \emph{raishū hatarakimasu} — \textquotedblleft{}I will work next week.\textquotedblright{}
\end{quote}

\subsection*{Your turn}
\begin{itemize}
\raggedright
  \item \emph{Say it:} \emph{raishū}
  \item \emph{Say it:} \emph{raishū hatarakimasu} — I will work next week
  \item \emph{Recall:} say \emph{last week}, \emph{this week}, \emph{next week}
\end{itemize}

\begin{tcolorbox}[breakable,colback=teal!4,colframe=teal!35!black,title={Before you move on}]
What is \textquotedblleft{}next week\textquotedblright{}? (\textbf{\ja{らいしゅう}}, \emph{raishū}.) Say it once more.
\end{tcolorbox}
\section[\texorpdfstring{\ja{らいげつ} (raigetsu)}{らいげつ (raigetsu)}]{\ja{らいげつ} (raigetsu) — next month}

\label{lesson:JA-C141-raigetsu}



\begin{quote}
Four recalls before the new word.

\begin{itemize}
\raggedright
  \item \emph{Recall:} say \emph{next week} — \textbf{R1}, one lesson back
  \item \emph{Recall:} say \emph{last month} — \textbf{R2}, five lessons back
  \item \emph{Recall:} say \emph{fluent} — \textbf{R3}, twenty lessons back
  \item \emph{Recall:} say \emph{this}, near the speaker — \textbf{R4}, eighty lessons back
\end{itemize}
\end{quote}

\subsection*{You'll want to know: \ja{らいげつ}}
\textbf{\ja{らいげつ}} — \emph{raigetsu} — \textquotedblleft{}next month\textquotedblright{}. The \textbf{\ja{らい}} at its start is the same as in \textbf{\ja{らいしゅう}}: \textquotedblleft{}the one coming\textquotedblright{}.

\begin{quote}
\textbf{\ja{らいげつ} \ja{うみに} \ja{いきます。}} — \emph{raigetsu umi ni ikimasu} — \textquotedblleft{}I will go to the sea next month.\textquotedblright{}
\end{quote}

\subsection*{Your turn}
\begin{itemize}
\raggedright
  \item \emph{Say it:} \emph{raigetsu}
  \item \emph{Say it:} \emph{raigetsu umi ni ikimasu} — I will go to the sea next month
  \item \emph{Recall:} say \emph{last month}, then \emph{next month}
\end{itemize}

\begin{tcolorbox}[breakable,colback=teal!4,colframe=teal!35!black,title={Before you move on}]
What is \textquotedblleft{}next month\textquotedblright{}? (\textbf{\ja{らいげつ}}, \emph{raigetsu}.) Say it once more.
\end{tcolorbox}
\section[\texorpdfstring{\ja{らいねん} (rainen)}{らいねん (rainen)}]{\ja{らいねん} (rainen) — next year}

\label{lesson:JA-C141-rainen}



\begin{quote}
Four recalls before the new word.

\begin{itemize}
\raggedright
  \item \emph{Recall:} say \emph{next month} — \textbf{R1}, one lesson back
  \item \emph{Recall:} say \emph{last year} — \textbf{R2}, five lessons back
  \item \emph{Recall:} write \textbf{\ja{ぺ}} — \textbf{R3}, twenty lessons back
  \item \emph{Recall:} write \textbf{\ja{れ}} — \textbf{R4}, eighty lessons back
\end{itemize}
\end{quote}

\subsection*{You'll want to know: \ja{らいねん}}
\textbf{\ja{らいねん}} — \emph{rainen} — \textquotedblleft{}next year\textquotedblright{}: \textbf{\ja{きょねん}}, \textbf{\ja{ことし}}, \textbf{\ja{らいねん}}, last year, this year, next year.

\begin{quote}
\textbf{\ja{らいねん} \ja{ともだちと} \ja{やまに} \ja{いきます。}} — \emph{rainen tomodachi to yama ni ikimasu} — \textquotedblleft{}Next year I will go to the mountains with a friend.\textquotedblright{} The particle \textbf{\ja{と}} means \textquotedblleft{}with\textquotedblright{}.
\end{quote}

\subsection*{Your turn}
\begin{itemize}
\raggedright
  \item \emph{Say it:} \emph{rainen}
  \item \emph{Say it:} \emph{rainen tomodachi to yama ni ikimasu} — Next year I will go to the mountains with a friend
  \item \emph{Recall:} say \emph{next week}, \emph{next month}, \emph{next year}
\end{itemize}

\begin{tcolorbox}[breakable,colback=teal!4,colframe=teal!35!black,title={Before you move on}]
What is \textquotedblleft{}next year\textquotedblright{}? (\textbf{\ja{らいねん}}, \emph{rainen}.) Say it once more.
\end{tcolorbox}
\section[\texorpdfstring{\ja{つもり} (tsumori)}{つもり (tsumori)}]{\ja{つもり} (tsumori) — an intention, a plan (I mean to ...)}

\label{lesson:JA-C141-tsumori}



\begin{quote}
Three recalls before the new word.

\begin{itemize}
\raggedright
  \item \emph{Recall:} say \emph{next year} — \textbf{R1}, one lesson back
  \item \emph{Recall:} say \emph{a moment ago} — \textbf{R2}, five lessons back
  \item \emph{Recall:} say \emph{that}, near the listener — \textbf{R4}, eighty lessons back
\end{itemize}
\end{quote}

\subsection*{You'll want to know: \ja{つもり}}
\textbf{\ja{つもり}} — \emph{tsumori} — \textquotedblleft{}an intention\textquotedblright{}, \textquotedblleft{}a plan\textquotedblright{}: what you mean to do. It follows a verb in its dictionary form, and \textbf{\ja{です}} closes the sentence.

\begin{quote}
\textbf{\ja{らいねん} \ja{うみに} \ja{いく} \ja{つもりです。}} — \emph{rainen umi ni iku tsumori desu} — \textquotedblleft{}I plan to go to the sea next year.\textquotedblright{}
\end{quote}

Here the dictionary form matters: \emph{iku}, not \emph{ikimasu}, comes before \textbf{\ja{つもり}}. \textbf{\ja{ます}} goes at the very end of a polite sentence, and \textbf{\ja{です}} does that job here.

\subsection*{Your turn}
\begin{itemize}
\raggedright
  \item \emph{Say it:} \emph{tsumori}
  \item \emph{Say it:} \emph{rainen umi ni iku tsumori desu} — I plan to go to the sea next year
  \item \emph{Recall:} say \emph{I plan to work next week}: \emph{raishū hataraku tsumori desu}
\end{itemize}

\begin{tcolorbox}[breakable,colback=teal!4,colframe=teal!35!black,title={Before you move on}]
What is \textquotedblleft{}a plan\textquotedblright{}, what you mean to do? (\textbf{\ja{つもり}}, \emph{tsumori}.) Say it once more.
\end{tcolorbox}
\section[\texorpdfstring{\ja{でしょう} (deshō)}{でしょう (deshō)}]{\ja{でしょう} (deshō) — probably; it will ..., I expect}

\label{lesson:JA-C141-deshou}



\begin{quote}
Four recalls before the new word.

\begin{itemize}
\raggedright
  \item \emph{Recall:} say \emph{a plan}, what you mean to do — \textbf{R1}, one lesson back
  \item \emph{Recall:} say \emph{it is over}, politely — \textbf{R2}, five lessons back
  \item \emph{Recall:} say \emph{I get up}, politely — \textbf{R3}, twenty lessons back
  \item \emph{Recall:} say \emph{a car} — \textbf{R4}, eighty lessons back
\end{itemize}
\end{quote}

\subsection*{You'll want to know: \ja{でしょう}}
\textbf{\ja{でしょう}} — \emph{deshō} — \textquotedblleft{}probably\textquotedblright{}, \textquotedblleft{}it will ..., I expect\textquotedblright{}. It closes a sentence about something you think will happen but cannot be sure of, the way a weather forecast does.

\begin{quote}
\textbf{\ja{あしたは} \ja{すずしいでしょう。}} — \emph{ashita wa suzushii deshō} — \textquotedblleft{}It will probably be cool tomorrow.\textquotedblright{}
\end{quote}

Now you can talk about what comes next: what will happen (\textbf{\ja{ます}} with a time word), what you mean to do (\textbf{\ja{つもり}}), and what will probably happen (\textbf{\ja{でしょう}}).

\subsection*{Your turn}
\begin{itemize}
\raggedright
  \item \emph{Say it:} \emph{deshō}
  \item \emph{Say it:} \emph{ashita wa suzushii deshō} — It will probably be cool tomorrow
  \item \emph{Recall:} say what comes next: \emph{raishū} (next week), \emph{raigetsu} (next month) and \emph{rainen} (next year), each with a verb in \emph{-masu}, then one plan with \emph{tsumori desu}, and one guess with \emph{deshō}
\end{itemize}

\begin{tcolorbox}[breakable,colback=teal!4,colframe=teal!35!black,title={Before you move on}]
Next week? (\textbf{\ja{らいしゅう}}, \emph{raishū}.) Next month? (\textbf{\ja{らいげつ}}, \emph{raigetsu}.) Next year? (\textbf{\ja{らいねん}}, \emph{rainen}.) A plan? (\textbf{\ja{つもり}}, \emph{tsumori}.) Probably? (\textbf{\ja{でしょう}}, \emph{deshō}.)
\end{tcolorbox}
